\section{Three Hecke factors strictly dominate the relative-zeta envelope}
\label{sec:s5-hecke-variational}

\paragraph{Review status.}
The complete post-writing author reading and two independent
whole-proof readings have passed. The source and finite-certificate
checks have their own exact recorded scope.

Keep the actual Calegari extension and the notation of MA-RB-002
and MA-RB-012. In particular, complex conjugation is of type
\(2B\), and the established four- and six-dimensional Artin
representations are modular. Put
\[
 \begin{aligned}
  R_B&=\eta\oplus W,&
  R_T&=W\oplus W\eta,&
  R_U&=V\eta\oplus X,\\
  L_B&=B,& L_T&=L_WL_{W\eta},& L_U&=L_{V\eta}L_X.
 \end{aligned}
 \tag{MA-RB.200}
\]
The subscript on \(R_B\) denotes a representation, not the
previous hypothesis \((\mathcal R_B)\).
Let \(N_j\) be the full Artin conductor of \(R_j\), for
\(j\in\{B,T,U\}\), and define
\[
 \begin{gathered}
 (d_j^+,d_j^-)=(4,2),(6,4),(4,6)
                         \quad(j=B,T,U),\\
 \Lambda_j(s)=N_j^{s/2}\Gamma_{\mathbb R}(s)^{d_j^+}
                \Gamma_{\mathbb R}(s+1)^{d_j^-}L_j(s),\\
 H_j(\sigma)=\frac{\Lambda_j'}{\Lambda_j}(\sigma)
   =\frac12\log N_j+d_j^+\psi_{\mathbb R}(\sigma)
       +d_j^-\psi_{\mathbb R}(\sigma+1)
       +\frac{L_j'}{L_j}(\sigma).
 \end{gathered}
 \tag{MA-RB.201}
\]
All local factors, including ramified factors, are retained.

\paragraph{MA-RB-013 (a strict three-factor refinement).}
The three functions in (MA-RB.200) are nonprincipal finite-order
Hecke \(L\)-functions over appropriate intermediate fields.
Their completions are entire of order at most one, nonzero at
zero and one, and have nonempty positive divisors in the
completed critical strip. Their real logarithmic derivatives
\(H_j\) extend continuously and positively to \([1,\infty)\).
For \(k\ge1\), set
\[
 \begin{gathered}
 F_{j,k}(\sigma)=(\sigma-\tfrac12)^2
                  -\frac{k(2\sigma-1)}{H_j(\sigma)},\\
 \mathfrak K_{L,k}=\min_{j\in\{B,T,U\}}
                            \inf_{\sigma>1}F_{j,k}(\sigma).
 \end{gathered}
 \tag{MA-RB.202}
\]
Each infimum is finite and is attained after adjoining
\(\sigma=1\). The following strict inequality is unconditional:
\[
             \mathfrak K_{L,k}<\mathfrak J_{L,k}<\frac14.
 \tag{MA-RB.203}
\]
Here \(\mathfrak J_{L,k}\) is the full finite subgroup envelope
of MA-RB-012, not only its quadratic member.

If comparison fails on the actual branch of MA-RB-002, its
real abscissa \(\gamma\) and integer pole order \(k\ge1\)
satisfy
\[
 \begin{gathered}
 \operatorname{ord}_{\gamma}\Lambda_j
    =\operatorname{ord}_{1-\gamma}\Lambda_j\ge k
                                  \quad(j=B,T,U),\\
 (\gamma-\tfrac12)^2\le\mathfrak K_{L,k}
                                     \le\mathfrak K_{L,1}.
 \end{gathered}
 \tag{MA-RB.204}
\]
Thus \(\mathfrak K_{L,1}\le0\) excludes that branch.
Otherwise its sufficient shared-zero gate is restricted to
\[
 B(\gamma)=L_U(\gamma)=0,\qquad
 \tfrac12<\gamma\le\tfrac12+\sqrt{\mathfrak K_{L,1}}<1.
 \tag{MA-RB.205}
\]
The additional necessary equality \(L_T(\gamma)=0\) also
holds. No assertion that this final interval is empty, or that
a positive-excess number field exists, is made. Neither
\(\mathsf A\), \(M\), nor \((\mathcal R_B)\) is used.

\begin{proof}
\emph{The actual induced Hecke factors.}
Let \(F_{20}\le S_5\) be the affine group on five letters.
The established permutation identity is
\(\operatorname{Ind}_{F_{20}}^{S_5}\mathbf1=\mathbf1\oplus W\eta\).
Its sign restriction is nontrivial, because \(F_{20}\) contains
a four-cycle. Twisting induction therefore gives
\[
       \operatorname{Ind}_{F_{20}}^{S_5}(\eta|_{F_{20}})
                      =\eta\oplus W=R_B.
 \tag{MA-RB.206}
\]
This is exactly the existing quadratic Hecke factor
\(\zeta_{EF}/\zeta_F\), where \(F=L^{F_{20}}\).

Embed \(A_4\) as the even permutations of four letters,
and let \(\chi\) be either nontrivial cubic character
of \(A_4/V_4\). Its values on the identity and the three
double transpositions are one, and its sums over the two
four-element classes of three-cycles are \(4\omega\) and
\(4\omega^2\). For an irreducible character \(\theta\)
of \(S_5\), Frobenius reciprocity consequently gives
\[
 \left\langle\theta,\operatorname{Ind}_{A_4}^{S_5}\chi\right\rangle
   =\frac{\theta(1)+3\theta(2B)-4\theta(3A)}{12}.
 \tag{MA-RB.207}
\]
The checked character table gives one for \(W,W\eta\)
and zero for \(\mathbf1,\eta,V,V\eta,X\). Hence
\(\operatorname{Ind}_{A_4}^{S_5}\chi=R_T\), for either
choice of the cubic character.

The action on oriented two-element subsets gives
\[
 \operatorname{Ind}_{S_3\times S_2}^{S_5}
                     (\mathbf1\boxtimes\operatorname{sgn}_{S_2})
      =\bigwedge\nolimits^2(\mathbf1\oplus V)=V\oplus X.
 \tag{MA-RB.208}
\]
Twisting by \(\eta\), whose restriction is the product
of the two signs, and using \(X\eta=X\), gives
\(\operatorname{Ind}_{S_3\times S_2}^{S_5}
(\operatorname{sgn}_{S_3}\boxtimes\mathbf1)=R_U\).
Each inducing character is nontrivial and finite order.
Artin induction, with the inertia-invariant local definition,
identifies each complete finite function with the corresponding
Hecke function. The exact supplied interface is
Hu--Kaneko--Martin--Schildkraut, \texttt{lem:artin-L}(a),(b),(d)
and \texttt{cor:hol-artin-induct}. These are classical inputs,
not new Hecke holomorphy assertions.

At complex conjugation the traces of \(R_B,R_T,R_U\)
are respectively \(2,2,-2\). Their dimensions are
\(6,10,10\), proving all the archimedean exponents in
(MA-RB.201). In particular, the exponent pair for \(T\)
is \((6,4)\): the sign character is trivial at this
complex conjugation.

\emph{Completions, endpoints and positive zero sums.}
There is also a direct Dedekind check of every completion.
Write
\[
 K_T=L^{A_4},\quad K_T'=L^{V_4},\qquad
 K_U=L^{S_3\times S_2},\quad K_U'=L^{A_3\times S_2}.
 \tag{MA-RB.209}
\]
The extensions \(K_T'/K_T\) and \(K_U'/K_U\) are cyclic
of degrees three and two. Artin induction and conductor
additivity, including the archimedean factors, give
\[
 \Lambda_B=\frac{\xi_{EF}}{\xi_F},\qquad
 \Lambda_T^2=\frac{\xi_{K_T'}}{\xi_{K_T}},\qquad
 \Lambda_U=\frac{\xi_{K_U'}}{\xi_{K_U}}.
 \tag{MA-RB.210}
\]
The normalization of \(\xi_K\) is exactly (MA-RB.181),
so both pole-removing factors \(s(s-1)\) cancel in each
ratio. In the cubic case the two conjugate characters
induce the same representation; this explains the square.
The corresponding conductor identities are
\(N_B=\Delta_{EF}/\Delta_F\),
\(N_T^2=\Delta_{K_T'}/\Delta_{K_T}\), and
\(N_U=\Delta_{K_U'}/\Delta_{K_U}\).

The precise classical Aramata input and the entire-completion
argument were checked in (MA-RB.177)--(MA-RB.185).
They apply to all three normal extensions in (MA-RB.210).
Thus the two unsquared completions are entire with positive
excess divisors, nonzero at both endpoints. The meromorphic
function \(\Lambda_T\) has an entire square, and is therefore
entire as well; its zero orders are half those of that square.
Its endpoint values are also nonzero. The canonical-product
order argument in MA-RB-012 applies to the ratios, and taking
a square root preserves the bound on order. Reflection and
real conjugation hold, with a possible constant sign in the
unsquared cubic completion; that sign disappears after
logarithmic differentiation.

The same already proved Hadamard-constant calculation gives
\[
 H_j(\sigma)=\sum_{\rho\in\mathcal Z_j}
          \frac{\sigma-\operatorname{Re}\rho}
               {|\sigma-\rho|^2},\qquad \sigma\ge1.
 \tag{MA-RB.211}
\]
Zeros are counted with their actual positive integer orders.
The sum converges absolutely and continuously at one.
The containing strip is \(0\le\operatorname{Re}\rho\le1\),
and the endpoints zero and one are absent. Nonemptiness
follows as in MA-RB-008/012: an empty divisor would force a
constant reflected order-one completion, whereas its logarithmic
derivative from (MA-RB.201) grows as
\(\tfrac12\dim R_j\log\sigma+O_L(1)\).
Therefore (MA-RB.211) is strictly positive on \([1,\infty)\).

Use the analogous notation \(H_Q=\xi_{\mathbb Q}'/\xi_{\mathbb Q}\),
\(H_F=\xi_F'/\xi_F\), and \(H_C,H_V,H_X\) for the completed
nontrivial factors \(C,L_V,L_X\). The latter two are entire
by Calegari's established four- and six-dimensional modularity;
their completed order-one functions and endpoint nonvanishing
have the standard cuspidal/Artin normalization. For \(V\),
one can check these facts directly from
\[
 K_5=L^{S_4},\qquad
 \Lambda_V=\frac{\xi_{K_5}}{\xi_{\mathbb Q}},\qquad
 N_V=\Delta_{K_5},\qquad (d_V^+,d_V^-)=(2,2).
 \tag{MA-RB.212}
\]
Here entireness of the finite \(L_V\) cancels all nontrivial
denominator zeros, and the two Dedekind endpoints are nonzero.
For \(X\), equivalently use the standard completed cuspidal
function furnished by the strong correspondence. No all-place
tensor-product identification is used. These completions also
have positive zero sums. The factor \(C\) is the usual
nontrivial quadratic Hecke function, and \(H_F\) is the
Dedekind positive sum. Consequently
\(H_C,H_V,H_X,H_F>0\) for \(\sigma\ge1\).

\emph{A quantitative comparison that pays for the trivial factor.}
We prove the following bound in precisely the required signature:
\[
                        H_V(\sigma)>H_Q(\sigma)
                                      \qquad(\sigma\ge1).
 \tag{MA-RB.213}
\]
No zero-location hypothesis on \(L_V\) is used.
First record a reflection-pair estimate. Put
\(x=\sigma-\tfrac12\), \(u=\operatorname{Re}\rho-\tfrac12\),
and \(y=(\operatorname{Im}\rho)^2\). A pair under
\(\rho\mapsto1-\overline\rho\), divided by \(2\sigma-1=2x\),
contributes
\[
 f(A)=\frac{A}{A^2+4u^2y},\qquad A=x^2+y-u^2.
 \tag{MA-RB.214}
\]
Fixed points are treated with half the paired weight.
For \(1\le\sigma\le3\), the endpoint values of \(A\) are
\(A_1=\tfrac14+y-u^2\) and \(A_3=\tfrac{25}4+y-u^2\).
The function \(f\) has no interior minimum, and
\[
 A_1A_3-4u^2y
    =y^2+(\tfrac{13}2-6u^2)y
          +(\tfrac14-u^2)(\tfrac{25}4-u^2)\ge0.
 \tag{MA-RB.215}
\]
Since \(u^2\le\tfrac14\), both factors are nonnegative;
the excluded completed endpoints remove the only singular
case. Thus \(f(A_1)\ge f(A_3)\), and summation gives
\[
 \frac{H_V(\sigma)}{2\sigma-1}\ge\frac{H_V(3)}5
                                      \quad(1\le\sigma\le3).
 \tag{MA-RB.216}
\]

For comparison, the elementary low-height bound needed for
\(\zeta_{\mathbb Q}\) follows from the classical split theta
integral. With \(\vartheta(x)=\sum_{n\ge1}e^{-\pi n^2x}\),
that formula in our normalization is
\[
 \xi_{\mathbb Q}(s)=1+s(s-1)
       \int_1^\infty\vartheta(x)
          \{x^{s/2}+x^{(1-s)/2}\}\,\frac{dx}{x}.
 \tag{MA-RB.217}
\]
The exact classical formulas are DLMF 20.10.2 and 20.7.32,
with the normalization DLMF 25.4.4; our \(\xi\) is twice
the latter. For \(0\le\operatorname{Re}s\le1\) and
\(|\operatorname{Im}s|\le1\), the absolute value of the
second term in (MA-RB.217) is bounded by
\[
 4\sum_{n\ge1}\frac{e^{-\pi n^2}}{\pi n^2}
       \le\frac{4e^{-\pi}}{\pi(1-e^{-3\pi})}<1.
 \tag{MA-RB.218}
\]
Here \(|s(s-1)|\le2\), and \(n^2\ge1+3(n-1)\).
Thus every completed zeta zero has \(|\operatorname{Im}\rho|>1\).
For its pair in (MA-RB.214), \(A\ge y>1\) when \(x\ge1/2\),
so \(A^2\ge4u^2y\). Hence \(f(A)\) decreases with \(x\), and
\[
 \frac{H_Q(\sigma)}{2\sigma-1}\le H_Q(1)
       =1+\frac{\gamma_{\!E}}2-\log2-\frac12\log\pi
       =:q\qquad(\sigma\ge1).
 \tag{MA-RB.219}
\]
The last equality follows by taking the limit at one in the
completed logarithmic derivative, using the Laurent constant
\(\gamma_{\!E}\) of zeta. This is the only low-height information
used, and (MA-RB.218) supplies it directly.

The quintic field \(K_5\) has signature \((1,2)\), because
\(2B\) fixes one letter. Minkowski's discriminant bound gives
\[
 \frac12\log N_V
    \ge\log\!\left((\pi/4)^2\frac{5^5}{5!}\right)
    =4\log5-\log3-7\log2+2\log\pi.
 \tag{MA-RB.220}
\]
The precise cited source is Milne (2020), version 3.08,
Theorem 4.3, printed p.~70, and the displayed discriminant
inequality in the proof of Theorem 4.9, printed p.~72.
Only this established Minkowski bound is used; it is not
reproved here.

Put \(S_3=\sum_{n\ge2}(\log n)n^{-3}<199/1000\).
At every ramified or unramified prime the local roots of
\(V\) have modulus one on the inertia invariants, of dimension
at most four. Thus, for real \(\sigma>1\),
\[
 E_V(\sigma):=-\frac{L_V'}{L_V}(\sigma)
       \le4E_Q(\sigma),\qquad
 E_Q(\sigma):=-\frac{\zeta'}\zeta(\sigma)
       \le\sum_{n\ge2}(\log n)n^{-\sigma}.
 \tag{MA-RB.221}
\]
The first expression is real; it need not be nonnegative.
The inequality is an upper bound on its actual local traces.
The gamma recurrence and special values (DLMF 5.4.12--5.4.13
and 5.5.2) give
\[
 \psi_{\mathbb R}(3)=1-\tfrac12\gamma_{\!E}-\log2-\tfrac12\log\pi,
 \quad
 \psi_{\mathbb R}(4)=\tfrac12-\tfrac12\gamma_{\!E}-\tfrac12\log\pi.
 \tag{MA-RB.222}
\]
Combining (MA-RB.220)--(MA-RB.222) yields
\[
 \begin{split}
 H_V(3)-5q\ \ge\;&4\log5-\log3-2
       -\tfrac92\gamma_{\!E}-4\log2+\tfrac52\log\pi-4S_3\\
 &>\frac3{125}>0.
 \end{split}
 \tag{MA-RB.223}
\]
Equations (MA-RB.216) and (MA-RB.219) prove (MA-RB.213)
on \([1,3]\).

For \(\sigma\ge3\), subtract the exact completed derivatives
and use (MA-RB.221):
\[
 \begin{split}
 H_V(\sigma)-H_Q(\sigma)\ge{}
 &\tfrac12\log N_V+\psi_{\mathbb R}(\sigma)
       +2\psi_{\mathbb R}(\sigma+1)\\
 &-\frac1\sigma-\frac1{\sigma-1}-3E_Q(\sigma).
 \end{split}
 \tag{MA-RB.224}
\]
The gamma derivatives increase, the two rational terms decrease,
and the positive Euler sum \(E_Q\) decreases. The right side is
therefore bounded below by its value with \(\sigma=3\) and
the same conductor lower bound. Explicitly it is greater than
\[
 4\log5-\log3+\tfrac76-\tfrac32\gamma_{\!E}
           -8\log2+\tfrac12\log\pi-3S_3
                     >\frac{179}{3000}>0.
 \tag{MA-RB.225}
\]
This finishes the proof of (MA-RB.213).

For clarity, the rational margins in the last two displays
use only
\[
 \log5>1.609,\quad \log3<1.099,\quad \log2<0.694,
 \quad\log\pi>1.144,\quad\gamma_{\!E}<0.578,
 \quad S_3<0.199.
 \tag{MA-RB.226}
\]
These bounds have an elementary finite certificate. For
\(1\le r\le2\), put \(z=(r-1)/(r+1)\) and use
\[
 2\sum_{j=0}^{15}\frac{z^{2j+1}}{2j+1}
 \ \le\ \log r\ \le\
 2\sum_{j=0}^{15}\frac{z^{2j+1}}{2j+1}
       +\frac{2z^{33}}{33(1-z^2)}.
 \tag{MA-RB.227}
\]
Reduce any positive rational argument by powers of two.
For \(\log\pi\), use \(\pi>157/50\) and apply this formula
to \(157/50\). The harmonic bound
\[
 \gamma_{\!E}\le\sum_{n=1}^{30}\frac1n-\log30-\frac1{62}
                    <0.578
 \tag{MA-RB.228}
\]
follows by summing
\(\log(1+1/n)-1/(n+1)\ge1/(2(n+1)^2)\).
Finally monotonicity of \((\log x)x^{-3}\) for \(x\ge2\)
gives
\[
 S_3\le\sum_{n=2}^{30}\frac{\log n}{n^3}
        +\frac{\tfrac12\log30+\tfrac14}{30^2}<0.199.
 \tag{MA-RB.229}
\]
Substitution of the rational upper bounds from (MA-RB.227)
verifies the last inequality. The companion rational-arithmetic
certificate records these inequalities; decimal floating-point
optimization and numerical zero searches are not used.

\emph{All subgroups, including the potentially adverse sign.}
Fix any nontrivial \(H\le S_5\), and set
\[
 \begin{gathered}
 c_\eta=1-h_\eta(H),\quad c_0=4-h_V(H),\quad
 c_1=4-h_{V\eta}(H),\quad c_X=6-h_X(H),\\
 a=5-h_W(H),\quad a'=5-h_{W\eta}(H),\qquad\delta=a-a'.
 \end{gathered}
 \tag{MA-RB.230}
\]
All coefficients are nonnegative integers. Since \(V\) is
faithful and \(H\ne1\), \(c_0\ge1\).
The character difference \(\chi_X-\chi_{V\eta}\), in the
class order \(1,2A,2B,3A,6A,4A,5A\), is
\((2,2,-2,-1,-1,0,2)\). It is everywhere at most two.
Averaging over \(H\) proves \(c_X\ge c_1\), and hence
\[
                            m_H=a+c_1.
 \tag{MA-RB.231}
\]
If \(H\subset A_5\), then \(\delta=0\).
Otherwise half its elements are odd, and the values of
\(\chi_W\) on the odd classes \(2A,6A,4A\) are
\(1,1,-1\). Therefore
\[
 \delta=h_{W\eta}(H)-h_W(H)
       =-\frac2{|H|}\sum_{h\in H\setminus A_5}\chi_W(h)
                   \in\{-1,0,1\}.
 \tag{MA-RB.232}
\]
In the two nonzero cases \(c_\eta=1\).

The genuine relative character is
\(R_H=\operatorname{Reg}-\operatorname{Ind}_H^{S_5}\mathbf1\).
If \(\delta\ge0\), direct coefficient comparison gives
\[
 R_H=a'R_T+\delta R_B+c_1R_U+c_0V
                    +(c_X-c_1)X+(c_\eta-\delta)\eta.
 \tag{MA-RB.233}
\]
Every coefficient is nonnegative. Taking the exact completed
logarithmic derivatives and using \(c_0\ge1\) gives
\[
 \mathscr H_H>a'H_T+\delta H_B+c_1H_U
                  \ge m_H\min\{H_B,H_T,H_U\}.
 \tag{MA-RB.234}
\]
This is an equality of all-place Artin factors before the
strict inequality; conductor and gamma terms are additive.

If \(\delta=-1\), the exact alternative is
\[
 R_H=aR_T+c_1R_U+c_0V+(c_X-c_1)X+\eta+W\eta.
 \tag{MA-RB.235}
\]
The completed meromorphic factor of \(W\eta\) is
\(\xi_F/\xi_{\mathbb Q}\), by the established permutation
identity. Its logarithmic derivative is \(H_F-H_Q\).
Consequently the leftover after \(aH_T+c_1H_U\) is
\[
 c_0H_V+(c_X-c_1)H_X+H_C+H_F-H_Q>0
 \tag{MA-RB.236}
\]
by (MA-RB.213) and \(c_0\ge1\).
No positivity of the meromorphic \(W\eta\) divisor is
asserted. Thus (MA-RB.234) holds in this case as well,
with weights \(a,c_1\) whose sum is again \(m_H\).

For every \(\sigma\ge1\), choose a minimizing member of
\(\{H_B,H_T,H_U\}\). Equation (MA-RB.234) implies
\[
       \min_j F_{j,k}(\sigma)<F_{H,k}(\sigma)
                                    \qquad(H\ne1).
 \tag{MA-RB.237}
\]
MA-RB-012 supplies a minimizing pair \((H_*,\sigma_*)\)
for its finite envelope after adjoining \(\sigma_*=1\).
Evaluate (MA-RB.237) there. Continuity allows this endpoint
value in the infimum over \(\sigma>1\). Hence
\(\mathfrak K_{L,k}<\mathfrak J_{L,k}\), proving (MA-RB.203).
The new minima themselves exist: the denominators are positive
and continuous at one, and (MA-RB.201) gives
\(H_j(\sigma)=\tfrac12\dim R_j\log\sigma+O_L(1)\)
as \(\sigma\to\infty\). Therefore \(F_{j,k}\to+\infty\).

\emph{Transfer of the actual comparison orders.}
At the real point \(\gamma\) in MA-RB-002, retain its notation
\(n,b,v_1,e,k\) and
\(x_1=\operatorname{ord}_\gamma L_{W\eta}\ge0\).
The latter nonnegativity is only at this real point and follows
from \(L_{W\eta}=\zeta_F/\zeta_{\mathbb Q}\), since the
denominator is nonzero on \((0,1)\). The established signed
orders give, without any additional nonvanishing premise,
\[
 \operatorname{ord}_\gamma L_B=b+n+k\ge k,\quad
 \operatorname{ord}_\gamma L_T=n+k+x_1\ge k,\quad
 \operatorname{ord}_\gamma L_U=v_1+e\ge k.
 \tag{MA-RB.238}
\]
All gamma factors are units at \(\gamma\), so these are
the completed orders. Reflection gives the same orders at
\(1-\gamma\). Keeping these two real zeros in (MA-RB.211)
gives
\[
 H_j(\sigma)\ge
     \frac{k(2\sigma-1)}
          {(\sigma-\tfrac12)^2-(\gamma-\tfrac12)^2}.
 \tag{MA-RB.239}
\]
Rearranging and taking the three infima proves (MA-RB.204).
Monotonicity in \(k\) proves its last inequality and the
gate statements.
\end{proof}

\paragraph{Exact finite evaluation and scope.}
For the three representations the actual ramified trace bound
is \(|\Lambda_{R_j}(n)|\le\dim(R_j)\log n\).
At a fixed \(\sigma>1\), truncating the Euler logarithmic
derivative at an integer \(N\ge3\) therefore has absolute error at most
\[
 \dim(R_j)N^{1-\sigma}
        \left(\frac{\log N}{\sigma-1}
                          +\frac1{(\sigma-1)^2}\right).
 \tag{MA-RB.240}
\]
Adding this error to the computed denominator gives a positive
upper denominator and hence a certified upper value of the
variational bound. Rational \(\sigma>1\) and increasing finite
truncations recover each infimum from above. This does not
assert uniform Euler convergence at \(\sigma=1\).

The representations in (MA-RB.206)--(MA-RB.208) are finite
Artin objects: they are induced rank-one local systems on
the corresponding finite covers. The exact finite-cover and
induction conventions are those already cited from SGA~1,
Expos\'e~V, \S7, and Deligne (1979), Definition~6.1 and
\S\S6.2--6.3. Conductor and gamma additivity in the proof
refer to these actual representations; no hypothetical higher
functorial transfer is introduced. The strict inequality is an
analytic comparison of their completed positive divisors.
It proves a smaller existing arithmetic gate, not equality
of the original five-dimensional functions.

\paragraph{Source and check boundary.}
Calegari's Theorem~1.1 supplies only its established
four- and six-dimensional strong correspondence here.
The new induced identities are explicitly verified above.
The classical Hecke, induction, Minkowski and theta formulas
are cited in their exact used forms, without replacing them
by new purported proofs. The new pair inequalities, numerical
margins, all-subgroup comparison and order transfer have
their own complete written-proof check events.
